% manuscript.tex — Finite-size scaling and percolation of cancer-cell-induced
%                  fluidisation in a phase field model of tissue
% Target: Physical Review E  (revtex4-2, two-column)
%
% Build:  pdflatex manuscript && pdflatex manuscript
%
% ============================================================================
% DRAFT CHECKLIST — mark [x] when complete, [~] when in progress
% ============================================================================
%
% --- 1. Validation (Palmieri Reproduction) ---
% [x] 1.1  Phase field snapshot at rho=0.90 (100 cells)
% [x] 1.2  L_n time series: soft=1.15, normal=1.06
% [x] 1.3  Displacement speed: soft 7.6% faster (p<0.001)
% [x] 1.4  Velocity distribution G(v_i), sigma_G=0.003
% [x] 1.5  MSD/4t plateau at 8tau
% [x] 1.6  D_eff extraction at lag=8tau validated
% [ ] 1.7  Fig 2-style deformation snapshot
% [ ] 1.8  Fig 4-style velocity distribution with Eq.5 fit
%
% --- 2. Finite-Size Scaling (Phase 1) ---
% [x] 2.1  N=100 rho=0.90: D_eff ratio = 1.34 +/- 0.13 (10 seeds, 200tau)
% [~] 2.2  N=400 rho=0.90: D_eff ratio = 1.25 +/- 0.17 (10 seeds, 67tau)
% [~] 2.3  N=200 rho=0.90: submitted, ~24h
% [ ] 2.4  N=800 rho=0.90: pending
% [ ] 2.5  N=1600 rho=0.90: pending (~5 days)
% [ ] 2.6  D_eff ratio vs N plot with error bars
% [ ] 2.7  L_n gap vs N plot
% [ ] 2.8  Mean speed ratio vs N plot
% [ ] 2.9  1/sqrt(N) extrapolation
% [ ] 2.10 rho=0.85 replicates
%
% --- 3. Percolation (Phase 3) ---
% [~] 3A.1-6  Pairwise: d=2R,4R,8R,20R running; 6R,12R pending
% [ ] 3A.7    Pairwise analysis: D_eff vs separation
% [~] 3B.1    f_c=0.05: 2/5 seeds running (chain 1 of 3)
% [ ] 3B.2-6  f_c=0.10-0.50: pending
% [ ] 3B.7    All-normal control: pending
% [ ] 3B.8-10 Percolation analysis and plots
% [ ] 3C.1-4  Clustered vs dispersed: pending
%
% --- 4. Figures ---
% [ ] 4.1  Fig 1: Model schematic + snapshot
% [ ] 4.2  Fig 2: Validation (L_n, velocity, MSD)
% [ ] 4.3  Fig 3: FSS D_eff ratio vs N
% [ ] 4.4  Fig 4: FSS L_n and speed vs N
% [ ] 4.5  Fig 5: Pairwise cooperativity
% [ ] 4.6  Fig 6: Percolation D_eff vs f_c
% [ ] 4.7  Fig 7: Clustered vs dispersed
%
% --- 5. Writing ---
% [~] 5.1  Introduction
% [~] 5.2  Model
% [ ] 5.3  Methods
% [ ] 5.4  Results: Validation
% [ ] 5.5  Results: FSS
% [ ] 5.6  Results: Percolation
% [ ] 5.7  Discussion
%
% --- 6. Infrastructure ---
% [x] 6.1  Rust study pipeline + TOML configs
% [x] 6.2  D_eff at 8tau in Rust
% [x] 6.3  Paired comparison with error propagation
% [x] 6.4  SVG plots from Rust
% [ ] 6.5  Cross-compile Rust for cluster
% [x] 6.6  fss.toml, percolation.toml
% [ ] 6.7  pairwise.toml
% ============================================================================
%
\documentclass[aps,pre,twocolumn,superscriptaddress,longbibliography,floatfix]{revtex4-2}

% ── packages ──────────────────────────────────────────────────────────────────
\usepackage{amsmath,amssymb}
\usepackage{graphicx}
\usepackage{bm}              % bold math
\usepackage{xcolor}
\usepackage{hyperref}
\usepackage{booktabs}        % nice tables
\usepackage{lineno}           % line numbers for review
\linenumbers

% ── convenience macros ────────────────────────────────────────────────────────
\newcommand{\dd}{\mathrm{d}}               % upright differential
\newcommand{\vA}{v_{\!A}}                  % self-propulsion speed
\newcommand{\peff}{p_{\mathrm{eff}}}       % effective shape index
\newcommand{\Deff}{D_{\mathrm{eff}}}       % effective diffusion coefficient
\newcommand{\phii}{\phi_i}                 % cell field shorthand
\newcommand{\phij}{\phi_j}                 % neighbor field shorthand
\newcommand{\fc}{f_c}                      % cancer cell fraction
\newcommand{\TODO}[1]{{\color{red}\textbf{[TODO: #1]}}}

% ══════════════════════════════════════════════════════════════════════════════
\begin{document}

\title{Finite-size scaling and percolation of cancer-cell-induced
fluidisation in a multi-cell phase field model}

\author{Steven~A.~Silber}
\affiliation{Department of Physics and Astronomy, Western University,
1151 Richmond Street, London, Ontario, Canada N6A 3K7}

\author{Mikko~Karttunen}
\email{mikko.karttunen1@uef.fi}
\affiliation{Department of Physics and Astronomy, Western University,
1151 Richmond Street, London, Ontario, Canada N6A 3K7}
\affiliation{ELLIS Institute Finland, Maarintie 8, 02150 Espoo, Finland}
\affiliation{Department of Technical Physics, University of Eastern
Finland, P.O.\ Box 1627, FI-70211 Kuopio, Finland}
\affiliation{Department of Chemistry, Western University,
1151 Richmond Street, London, Ontario, Canada N6A 5B7}

\date{\today}

% ── abstract ──────────────────────────────────────────────────────────────────
\begin{abstract}
Palmieri \textit{et al.}\ showed that a single soft cell embedded in a
stiff monolayer of 72 phase field cells exhibits enhanced motility
through intermittent speed bursts driven by elastic mismatch.  Whether
this result survives in larger systems, and what happens when multiple
soft cells are present, has not been tested.  We perform a systematic
finite-size study of the Palmieri model with $N = 100$ to $12{,}800$
cells at packing fractions $\rho = 0.70$--$1.00$ and find that the
motility enhancement ratio $\Deff^{\mathrm{cancer}} /
\Deff^{\mathrm{normal}}$ \TODO{converges / diverges / vanishes} as
$N \to \infty$.  For systems with a cancer cell fraction~$\fc$, the
fraction of mobile cells undergoes a percolation transition at
$\fc^* = $ \TODO{value}, with the cluster size distribution consistent
with the 2D percolation universality class ($\tau_p \approx 2.05$).
Below~$\fc^*$, mobile regions are isolated around individual cancer
cells; above it, a system-spanning mobile network forms.  The
fluidisation range of each cancer cell, measured by the excess
diffusion coefficient as a function of distance, decays
\TODO{exponentially / algebraically} with characteristic length
$\xi = $ \TODO{value}.  These results place quantitative bounds on the
system sizes required for convergence and on the cancer cell fraction
needed for collective tissue fluidisation.
\end{abstract}

\maketitle

% ══════════════════════════════════════════════════════════════════════════════
% TEMPORARY: Draft checklist (remove before submission)
% ══════════════════════════════════════════════════════════════════════════════
\onecolumngrid
\begingroup
\small
\color{red}
\section*{DRAFT CHECKLIST (remove before submission)}

\subsection*{1. Validation (Palmieri Reproduction)}
\begin{itemize}
\item[$\checkmark$] 1.1~~Phase field snapshot at $\rho=0.90$ (100 cells)
\item[$\checkmark$] 1.2~~$L_n$ time series: soft $=1.15$, normal $=1.06$
\item[$\checkmark$] 1.3~~Displacement speed: soft 7.6\% faster ($p<0.001$)
\item[$\checkmark$] 1.4~~Velocity distribution $G(v_i)$, $\sigma_G=0.003$
\item[$\checkmark$] 1.5~~MSD$/4t$ plateau at $8\tau$
\item[$\checkmark$] 1.6~~$\Deff$ extraction at lag $= 8\tau$ validated
\item[$\square$] 1.7~~Fig~2-style deformation snapshot
\item[$\square$] 1.8~~Fig~4-style velocity distribution with Eq.~5 fit
\end{itemize}

\subsection*{2. Finite-Size Scaling}
\begin{itemize}
\item[$\checkmark$] 2.1~~$N=100$, $\rho=0.90$: $\Deff$ ratio $= 1.34 \pm 0.13$ (10 seeds, $200\tau$)
\item[$\sim$] 2.2~~$N=400$, $\rho=0.90$: $\Deff$ ratio $= 1.25 \pm 0.17$ (10 seeds, $67\tau$ prelim.)
\item[$\sim$] 2.3~~$N=200$, $\rho=0.90$: submitted, $\sim$24h
\item[$\square$] 2.4~~$N=800$: pending
\item[$\square$] 2.5~~$N=1600$: pending ($\sim$5 days)
\item[$\square$] 2.6~~$\Deff$ ratio vs $N$ plot with error bars
\item[$\square$] 2.7~~$L_n$ gap vs $N$ plot
\item[$\square$] 2.8~~Mean speed ratio vs $N$ plot
\item[$\square$] 2.9~~$1/\sqrt{N}$ extrapolation
\item[$\square$] 2.10~~$\rho=0.85$ replicates
\end{itemize}

\subsection*{3. Percolation}
\begin{itemize}
\item[$\sim$] 3A~~Pairwise: $d=2R,4R,8R,20R$ running; $6R,12R$ pending
\item[$\square$] 3A.7~~Analysis: $\Deff$ vs separation
\item[$\sim$] 3B.1~~$\fc=0.05$: 2/5 seeds running (chain 1/3)
\item[$\square$] 3B.2--6~~$\fc=0.10$--$0.50$: pending
\item[$\square$] 3B.8--10~~Percolation analysis and plots
\item[$\square$] 3C~~Clustered vs dispersed: pending
\end{itemize}

\subsection*{4. Figures}
\begin{itemize}
\item[$\square$] Fig~1: Model schematic + snapshot
\item[$\square$] Fig~2: Validation ($L_n$, velocity, MSD)
\item[$\square$] Fig~3: FSS $\Deff$ ratio vs $N$
\item[$\square$] Fig~4: FSS $L_n$ and speed vs $N$
\item[$\square$] Fig~5: Pairwise cooperativity
\item[$\square$] Fig~6: Percolation $\Deff$ vs $\fc$
\item[$\square$] Fig~7: Clustered vs dispersed
\end{itemize}

\subsection*{5. Writing}
\begin{itemize}
\item[$\sim$] Introduction ~~\item[$\sim$] Model ~~\item[$\square$] Methods
\item[$\square$] Results: Validation ~~\item[$\square$] Results: FSS
\item[$\square$] Results: Percolation ~~\item[$\square$] Discussion
\end{itemize}
\endgroup
\twocolumngrid

% ══════════════════════════════════════════════════════════════════════════════
%  I.  INTRODUCTION
% ══════════════════════════════════════════════════════════════════════════════
\section{Introduction}
\label{sec:intro}

The mechanical interplay between cancer cells and surrounding healthy
tissue governs several stages of tumour progression.  Cancer cells in
dense epithelia must navigate a jammed environment of tightly packed
neighbours, and the degree to which a mechanically distinct cell can
fluidise its local environment determines its capacity for invasion
and metastasis~\cite{Park2015,Mongera2018,Grosser2021}.

Palmieri \textit{et al.}~\cite{Palmieri2015} demonstrated this effect
in a multi-cell phase field model: a single cell with reduced
stiffness ($\gamma_c / \gamma_n = 0.35$), embedded in a monolayer of
72 cells at packing fraction $\rho = 0.85$--$0.90$, exhibits
intermittent speed bursts, non-Gaussian velocity distributions, and an
effective diffusion coefficient $\Deff$ elevated above the all-normal
baseline.  The mechanism is elastic mismatch: the soft cell deforms
more easily through the cage of stiff neighbours.  Lee
\textit{et al.}~\cite{Lee2012} observed qualitatively similar
intermittent migration of MDA-MB-231 breast cancer cells in confluent
MCF-10A monolayers.

The Palmieri result raises three questions that the original 72-cell
study could not address.  First, 72 cells ($\sim 8$ cells per linear
dimension) is a small system.  The cancer cell's cage has only $\sim
6$ neighbours, and the periodic images are separated by fewer than 4
cell diameters.  Whether the motility enhancement, the burst
statistics, and the velocity distribution parameters converge at this
system size is unknown.  Second, real tumours contain populations of
mechanically distinct cells, not isolated individuals.  When the
cancer cell fraction~$\fc$ increases, the mobile regions around
individual cancer cells may connect into a system-spanning network,
producing a percolation transition from local to global
fluidisation~\cite{Angelini2011}.  Third, monodisperse simulations
at high packing fraction develop hexagonal ordering artefacts that
suppress the amorphous packing found in real tissues.

We address the first two of these questions.  Using the same phase
field model as Palmieri \textit{et al.}, we scale the system from
$N = 100$ to $12{,}800$ cells, sweep the packing fraction from
$\rho = 0.70$ (non-confluent) to $1.00$ (fully confluent), and
increase the cancer cell fraction from a single cell ($\fc = 1/N$)
to $\fc = 0.50$.

The model is defined in Sec.~\ref{sec:model} and simulation methods
in Sec.~\ref{sec:methods}.


% ══════════════════════════════════════════════════════════════════════════════
%  II.  MODEL
% ══════════════════════════════════════════════════════════════════════════════
\section{Model}
\label{sec:model}

We use the multi-cell phase field model of Palmieri
\textit{et al.}~\cite{Palmieri2015} without modification.  Each of
$N$ cells in a two-dimensional square domain of side~$L$ with periodic
boundary conditions is described by a scalar order parameter
$\phii(\bm{r},t) \in [0,1]$, with $\phii \approx 1$ inside cell~$i$
and $\phii \approx 0$ outside.  The total free energy is
%
\begin{equation}
  F = F_{\mathrm{int}} + F_{\mathrm{vol}} + F_{\mathrm{rep}},
  \label{eq:free-energy}
\end{equation}
%
where the single-cell term contains a gradient energy and double-well
potential,
%
\begin{equation}
  F_{\mathrm{int}} = \sum_i \int \biggl[
    \gamma_i |\nabla\phii|^2
    + \frac{30\gamma_i}{\lambda^2}\,\phii^2(1-\phii)^2
  \biggr] \dd A,
  \label{eq:F-int}
\end{equation}
%
a soft volume constraint penalises deviations from the target area
$A_0 = \pi R^2$,
%
\begin{equation}
  F_{\mathrm{vol}} = \mu \sum_i
    \biggl(\int \phii^2 \, \dd A - A_0\biggr)^{\!2},
  \label{eq:F-vol}
\end{equation}
%
and quartic repulsion prevents cell overlap,
%
\begin{equation}
  F_{\mathrm{rep}} = \kappa \sum_{i<j} \int \phii^2\,\phij^2 \; \dd A.
  \label{eq:F-rep}
\end{equation}
%
The gradient energy coefficient~$\gamma_i$ is the per-cell stiffness
parameter that distinguishes cancer cells ($\gamma_c$) from normal
cells ($\gamma_n$).  All other parameters ($\kappa$, $\mu$, $\lambda$,
$R$) are identical across cell types.

The equation of motion for each cell field
is~\cite{Palmieri2015}
%
\begin{equation}
  \frac{\partial \phii}{\partial t}
    + \bm{v}_i \cdot \nabla\phii
    = -M\,\frac{\delta F}{\delta \phii},
  \label{eq:eom}
\end{equation}
%
where $M = 1/2$ is the mobility.  The cell velocity
$\bm{v}_i = \bm{v}_{i,I} + \bm{v}_{i,A}$ decomposes into
an interaction velocity $\bm{v}_{i,I}$ derived from the
dissipation principle of Palmieri \textit{et al.}\ (parameterised by
the friction coefficient~$\xi$) and an active self-propulsion
$\bm{v}_{i,A} = \vA\,\hat{\bm{p}}_i$.  The polarity
$\hat{\bm{p}}_i$ reorients via a run-and-tumble process with
persistence time~$\tau$.

%% Cancer cell definition
A cancer cell is defined solely by a reduced stiffness: $\gamma_c <
\gamma_n$.  Following Palmieri \textit{et al.}, the ratio
$\gamma_c / \gamma_n = 0.35$ is used throughout unless stated
otherwise.  In runs with a cancer cell fraction~$\fc$, cells are
randomly assigned as cancer with probability~$\fc$ at initialisation.


% ══════════════════════════════════════════════════════════════════════════════
%  III.  SIMULATION METHODS
% ══════════════════════════════════════════════════════════════════════════════
\section{Simulation methods}
\label{sec:methods}

\subsection{Implementation}
\label{sec:methods:gpu}

The model is implemented as a custom GPU-accelerated (CUDA C++) code.
All $N$ phase fields are stored on a single GPU and evolved
simultaneously using per-cell bounding-box subdomains (side
$\approx 2.5 \times 2R$) rather than the full $L \times L$ grid,
reducing GPU memory from $O(NL^2)$ to $O(NR^2)$ and enabling
simulations up to $N = 12{,}800$ on a single NVIDIA H100 GPU.

\subsection{Parameters}
\label{sec:methods:params}

Table~\ref{tab:params} lists the simulation parameters, which are
identical to those of Palmieri
\textit{et al.}~\cite{Palmieri2015}.

\begin{table}[b]
  \caption{Simulation parameters, identical to Palmieri
    \textit{et al.}~\cite{Palmieri2015}.  The stiffness
    ratio $\gamma_c / \gamma_n = 0.35$ defines the cancer cell.}
  \label{tab:params}
  \begin{ruledtabular}
  \begin{tabular}{lll}
    Parameter & Symbol & Value \\ \hline
    Cell radius        & $R$        & 49                \\
    Interface width    & $\lambda$  & 7                 \\
    Normal stiffness   & $\gamma_n$ & 1.0               \\
    Cancer stiffness   & $\gamma_c$ & 0.35              \\
    Repulsion          & $\kappa$   & 10                \\
    Volume penalty     & $\mu$      & 1.0               \\
    Friction           & $\xi$      & 1500              \\
    Mobility           & $M$        & 0.5               \\
    Active speed       & $\vA$      & 0.01              \\
    Tumble time        & $\tau$     & $10{,}000$        \\
    Time step          & $\Delta t$ & 0.02              \\
    Boundary cond.     &            & periodic          \\
  \end{tabular}
  \end{ruledtabular}
\end{table}

The packing fraction $\rho = N\pi R^2 / L^2$ is varied from 0.70 to
1.00 by adjusting the domain side~$L$.  The system size~$N$ ranges
from 100 to 12{,}800 cells.

\subsection{Equilibration}
\label{sec:methods:eq}

For each $(N, \rho)$ pair, we equilibrate an all-normal ($\gamma_i
= \gamma_n$ for all~$i$) configuration at $\vA = 0$ for
$t = 80{,}000$ ($8\tau$), following the protocol of Palmieri
\textit{et al.}~\cite{Palmieri2015}.  Cells are initialised on a
hexagonal grid and relaxed until all transient stress has decayed.
The resulting checkpoints serve as independent starting states for
production runs.  We verify convergence by confirming that control
runs produce mean centroid displacements negligible compared to~$R$.

\subsection{Production protocol}
\label{sec:methods:production}

\paragraph{Phase 1: Finite-size scaling (single cancer cell).}
%
Starting from an equilibrated checkpoint, we set cell~0 to $\gamma_c
= 0.35$ and enable motility $\vA = 0.01$ for all cells.  Production
runs last $t = 100{,}000$ ($10\tau$).  An all-normal control is run
from the same checkpoint.

\paragraph{Phase 3: Multiple cancer cells.}
%
Starting from an equilibrated checkpoint, a fraction~$\fc$ of cells
is randomly assigned $\gamma_c = 0.35$.  The same motility,
duration, and control protocol apply.

\subsection{Observables}
\label{sec:methods:obs}

\paragraph{Effective diffusion coefficient.}
The long-time diffusion coefficient is extracted from the
mean-squared displacement,
\begin{equation}
  \Deff = \lim_{t\to\infty} \frac{\langle\Delta r^2(t)\rangle}{4t},
  \label{eq:Deff}
\end{equation}
computed separately for cancer cells, normal cells, and the full
population.

\paragraph{Burst statistics.}
Following Palmieri \textit{et al.}~\cite{Palmieri2015}, a speed
burst is defined as an event where $|v_i| > \mu_v + 3\sigma_v$
(where $\mu_v$ and $\sigma_v$ are the tagged cell's mean and standard
deviation of speed) for at least $T_{\mathrm{burst}}$ consecutive
time steps.  We report the burst frequency $f_b$, mean duration, and
mean peak amplitude.

\paragraph{Velocity distributions.}
The velocity components $v_x$ and $v_y$ of the cancer cell are
histogrammed and fit to a Student-$t$ distribution with
$\nu$~degrees of freedom.  Heavy tails ($\nu < 5$) indicate
intermittent dynamics.

\paragraph{Self-overlap function.}
\begin{equation}
  Q(\Delta t) = \frac{1}{N}\sum_{i=1}^{N}
  \Theta\!\bigl(a - |\bm{r}_i(t_0 + \Delta t) - \bm{r}_i(t_0)|\bigr),
  \label{eq:Qt}
\end{equation}
where $a = 0.3 \times 2R$.  We fit $Q(t) = \exp[-(t/\tau_\alpha)^\beta]$
to extract the structural relaxation time $\tau_\alpha$ and stretching
exponent~$\beta$.

\paragraph{Effective shape index.}
From the per-cell perimeter $L_i$ recorded in the trajectory,
\begin{equation}
  \peff = L_i \times 2\sqrt{\pi},
  \label{eq:peff}
\end{equation}
which equals $P_i / \sqrt{A_i}$ for a cell whose area is constrained
to~$A_0$.

\paragraph{Percolation observables (Phase 3).}
Mobile cells are identified as those whose cage-relative displacement
exceeds a threshold over a time window $T_{\mathrm{obs}}$.  Connected
clusters of mobile cells are determined via Voronoi adjacency.  The
cluster size distribution $P(s)$, the order parameter
$P_\infty = s_{\mathrm{max}} / N$ (largest cluster fraction), and the
susceptibility $\chi = \langle s^2\rangle / \langle s\rangle$
(excluding the largest cluster) characterise the percolation
transition.

\paragraph{Mobility profile.}
The diffusion coefficient of normal cells is computed as a function
of their distance from the nearest cancer cell, $\Deff(r)$.  This
quantifies the spatial range of the fluidisation induced by a single
cancer cell.


% ══════════════════════════════════════════════════════════════════════════════
%  IV.  RESULTS — Finite-size scaling
% ══════════════════════════════════════════════════════════════════════════════
\section{Finite-size scaling}
\label{sec:fss}

\subsection{Motility enhancement ratio}
\label{sec:fss:ratio}

\TODO{$\Deff^{\mathrm{cancer}} / \Deff^{\mathrm{normal}}$ vs $N$ at
$\rho = 0.85$ and $0.90$.  Does it converge, diverge, or vanish?
Include error bars from 10--20 realisations.}

\subsection{Burst statistics}
\label{sec:fss:bursts}

\TODO{Burst frequency, duration, and amplitude vs $N$.  Do they
converge?  At what system size?}

\subsection{Velocity distributions}
\label{sec:fss:velocity}

\TODO{Student-$t$ degrees of freedom $\nu$ vs $N$.  Does the
heavy-tailed character persist at large $N$?}

\subsection{Packing fraction dependence}
\label{sec:fss:rho}

\TODO{$\Deff$ ratio vs $\rho$ at the converged system size.
Map the $(N, \rho)$ space.  At what $\rho$ does the enhancement
disappear (cells too far apart for elastic mismatch to matter)?}


% ══════════════════════════════════════════════════════════════════════════════
%  V.  RESULTS — Multiple cancer cells and percolation
% ══════════════════════════════════════════════════════════════════════════════
\section{Cancer cell populations and percolation}
\label{sec:percolation}

\subsection{Effective diffusion vs cancer fraction}
\label{sec:percolation:Deff}

\TODO{$\Deff$ of cancer cells and normal cells vs $\fc$ at
$N = 4000$, $\rho = 0.90$.  Include the $\fc = 1/N$ (Phase~1)
data point as the single-cell limit.}

\subsection{Percolation transition}
\label{sec:percolation:transition}

\TODO{$P_\infty$, $\langle s \rangle$, $\chi$ vs $\fc$.  Locate
$\fc^*$ from the peak of~$\chi$.  Fit $P(s)$ at $\fc^*$ to a
power law $s^{-\tau_p}$ and compare to the 2D percolation exponent
$\tau_p = 187/91 \approx 2.05$.  Finite-size scaling with
$N = 2000$, $4000$, $6000$.}

\subsection{Mobility profile}
\label{sec:percolation:profile}

\TODO{$\Deff(r | \text{cancer})$ for normal cells as a function of
distance from the nearest cancer cell.  Extract the fluidisation
length~$\xi$ from an exponential or power-law fit.  Compare $\xi$
to the inter-cancer-cell distance at $\fc^*$ to verify that
percolation occurs when fluidisation zones overlap.}

\subsection{Pairwise cooperativity}
\label{sec:percolation:pair}

\TODO{Two cancer cells at separation $d$ in $N = 2000$.  Excess
$\Deff$ vs $d$.  Does the interaction decay exponentially or
algebraically?  What is the characteristic interaction length?}

\subsection{Clustered vs dispersed}
\label{sec:percolation:geometry}

\TODO{At fixed $\fc$, compare $\Deff$ for cancer cells in a compact
cluster vs randomly dispersed.  Prediction: interior cancer cells in
the cluster are less motile (they lack elastic mismatch contact with
surrounding stiff tissue).}


% ══════════════════════════════════════════════════════════════════════════════
%  VI.  DISCUSSION
% ══════════════════════════════════════════════════════════════════════════════
\section{Discussion}
\label{sec:discussion}

\TODO{Points to cover:

1.\ Whether the Palmieri single-cell result is a finite-size artefact
or survives to large $N$.  If $\Deff^{\mathrm{cancer}} /
\Deff^{\mathrm{normal}}$ converges, the mechanism is robust; if it
vanishes, the original result is a boundary-condition effect.

2.\ The percolation threshold $\fc^*$ and whether it falls in the 2D
percolation universality class.  Physical interpretation: at $\fc^*$,
the fluidisation zones of individual cancer cells overlap and form a
connected pathway.  The length scale $\xi$ sets the typical zone
radius; when $\xi \sim f_c^{-1/2}$ (mean inter-cancer distance),
percolation occurs.

3.\ Connection to vertex model literature.  The shape index
$\peff$ at $\fc^*$ can be compared to the vertex model transition
at $p_0^* \approx 3.81$~\cite{Bi2015}.  If $\langle\peff\rangle$
crosses 3.81 at the percolation threshold, the elastic-mismatch
fluidisation maps onto the same order parameter.

4.\ Packing fraction dependence.  At low $\rho$ (non-confluent),
cells are not in contact; elastic mismatch requires cell-cell contact
and therefore $\rho$ near confluence.  The $\rho$ dependence of the
enhancement ratio maps the boundary between the contact-dependent
(confluent) and gap-dominated (non-confluent) regimes.

5.\ Limitations: monodisperse cell sizes (polydispersity study
deferred), no adhesion, 2D only.  Extension to polydispersity will
suppress hexagonal ordering and test whether the motility threshold
is sharp or continuous.

6.\ Comparison with experiment: Lee \textit{et al.}~\cite{Lee2012}
observed intermittent migration of MDA-MB-231 cells in MCF-10A
monolayers.  The burst frequency and velocity distribution shape
can be compared quantitatively once finite-size convergence is
established.}


% ══════════════════════════════════════════════════════════════════════════════
%  ACKNOWLEDGMENTS
% ══════════════════════════════════════════════════════════════════════════════
\begin{acknowledgments}
Computational resources were provided by the Digital Research Alliance
of Canada.
M.K.\ thanks the Natural Sciences and Engineering Research Council of
Canada (NSERC), the Canada Research Chairs Program, and Foundation PS
for financial support.
S.A.S.\ thanks NSERC for financial support through the Postgraduate
Scholarships (PGS~D) program.
\end{acknowledgments}


% ══════════════════════════════════════════════════════════════════════════════
%  BIBLIOGRAPHY
% ══════════════════════════════════════════════════════════════════════════════
\begin{thebibliography}{99}

\bibitem{Palmieri2015}
B.~Palmieri, Y.~Bresler, D.~Wirtz, and M.~Grant,
Multiple scale model for cell migration in monolayers,
\href{https://doi.org/10.1038/srep11745}{Sci.\ Rep.\ \textbf{5}, 11745
(2015)}.

\bibitem{Lee2012}
M.~H. Lee, P.~H. Wu, J.~K. Strimbu, C.~N. Bhargava, and D.~Bhatt,
Mismatch in mechanical and adhesive properties induces
pulsating cancer cell migration in epithelial monolayer,
\href{https://doi.org/10.1016/j.bpj.2012.03.027}{Biophys.\ J.\
\textbf{102}, 2731 (2012)}.

\bibitem{Park2015}
J.~A. Park \textit{et~al.},
Unjamming and cell shape in the asthmatic airway epithelium,
\href{https://doi.org/10.1038/nmat4357}{Nat.\ Mater.\ \textbf{14}, 1040
(2015)}.

\bibitem{Mongera2018}
A.~Mongera \textit{et~al.},
A fluid-to-solid jamming transition underlies vertebrate body axis
elongation,
\href{https://doi.org/10.1038/s41586-018-0479-2}{Nature \textbf{561}, 401
(2018)}.

\bibitem{Grosser2021}
S.~Grosser \textit{et~al.},
Cell and nucleus shape as an indicator of tissue fluidity in carcinoma,
\href{https://doi.org/10.1103/PhysRevX.11.011033}{Phys.\ Rev.\ X
\textbf{11}, 011033 (2021)}.

\bibitem{Angelini2011}
T.~E. Angelini \textit{et~al.},
Glass-like dynamics of collective cell migration,
\href{https://doi.org/10.1073/pnas.1010059108}{Proc.\ Natl.\ Acad.\ Sci.\
USA \textbf{108}, 4714 (2011)}.

\bibitem{Bi2015}
D.~Bi, J.~H. Lopez, J.~M. Schwarz, and M.~L. Manning,
A density-independent glass transition in biological tissues,
\href{https://doi.org/10.1038/nphys3471}{Nat.\ Phys.\ \textbf{11}, 1074
(2015)}.

\bibitem{Bi2016}
D.~Bi, X.~Yang, M.~C. Marchetti, and M.~L. Manning,
Motility-driven glass and jamming transitions in biological tissues,
\href{https://doi.org/10.1103/PhysRevX.6.021011}{Phys.\ Rev.\ X
\textbf{6}, 021011 (2016)}.

\bibitem{NagaiHonda2001}
T.~Nagai and H.~Honda,
A dynamic cell model for the formation of epithelial tissues,
\href{https://doi.org/10.1080/13642810108205772}{Philos.\ Mag.\ B
\textbf{81}, 699 (2001)}.

\bibitem{Farhadifar2007}
R.~Farhadifar, J.-C. R\"oben, B.~Aigouy, S.~Eaton, and F.~J\"ulicher,
The influence of cell mechanics, cell-cell interactions, and
proliferation on epithelial packing,
\href{https://doi.org/10.1016/j.cub.2007.11.049}{Curr.\ Biol.\
\textbf{17}, 2095 (2007)}.

\bibitem{Atia2018}
L.~Atia \textit{et~al.},
Geometric constraints during epithelial jamming,
\href{https://doi.org/10.1038/s41567-018-0089-9}{Nat.\ Phys.\ \textbf{14},
613 (2018)}.

\bibitem{Camley2017}
B.~A. Camley and W.-J. Rappel,
Physical models of collective cell motility: from cell to tissue,
\href{https://doi.org/10.1088/1361-6463/aa56fe}{J.\ Phys.\ D \textbf{50},
113002 (2017)}.

\bibitem{Nonomura2012}
M.~Nonomura,
Study on multicellular systems using a phase field model,
\href{https://doi.org/10.1371/journal.pone.0033501}{PLoS ONE \textbf{7},
e33501 (2012)}.

\bibitem{Loewe2020}
B.~Loewe, M.~Chiang, D.~Marenduzzo, and M.~C. Marchetti,
Solid-liquid transition of deformable and overlapping active particles,
\href{https://doi.org/10.1103/PhysRevLett.125.038003}{Phys.\ Rev.\ Lett.\
\textbf{125}, 038003 (2020)}.

\bibitem{Saito2024}
S.~Saito and S.~Ishihara,
Cell deformability gives rise to a fluid-to-fluid transition in tissues,
\href{https://doi.org/10.1126/sciadv.adi8433}{Sci.\ Adv.\ \textbf{10},
eadi8433 (2024)}.

\bibitem{Garcia2015}
S.~Garcia, E.~Hannezo, J.~Elgeti, J.-F. Joanny, P.~Silberzan, and
N.~S. Gov,
Physics of active jamming during collective cellular motion in a
monolayer,
\href{https://doi.org/10.1073/pnas.1510973112}{Proc.\ Natl.\ Acad.\ Sci.\
USA \textbf{112}, 15314 (2015)}.

\bibitem{Malinverno2017}
C.~Malinverno \textit{et~al.},
Endocytic reawakening of motility in jammed epithelia,
\href{https://doi.org/10.1038/nmat4848}{Nat.\ Mater.\ \textbf{16}, 587
(2017)}.

\bibitem{Bresler2018}
Y.~Bresler, B.~Palmieri, and M.~Grant,
Sharp interface limit of the phase field model for cell motility
(2018), \href{https://arxiv.org/abs/1807.10318}{arXiv:1807.10318}.

\bibitem{HohenbergHalperin1977}
P.~C. Hohenberg and B.~I. Halperin,
Theory of dynamic critical phenomena,
\href{https://doi.org/10.1103/RevModPhys.49.435}{Rev.\ Mod.\ Phys.\
\textbf{49}, 435 (1977)}.

\end{thebibliography}

\end{document}
